\section{Implementierung des Attention-basierten Selektionsmodells}
\label{sec:implementierung_des_attention_basierten_selektionsmodells}

Die Gewichtungsfunktion wird für jeden CTI-Report separat berechnet, wodurch keine wechselseitigen Abhängigkeiten zwischen den einzelnen Reports entstehen. Dadurch eignen sich sowohl die Gewichtungsberechnung als auch die anschließende Aggregation grundsätzlich für eine parallele Verarbeitung über mehrere CTI-Reports hinweg.\\
Im Gegensatz zu frühen spektralen Ansätzen in der Graph-Verarbeitung sind keine rechenintensiven Matrixoperationen wie Inversionen oder Eigenwertzerlegungen erforderlich \cite{velickovic2018gat}.
% es gibt keine abhängigkeiten zwischen report A und B die berechnung wird für jeden report einzeln durchgeführt und wir müssen keine komplexe graphen matrix aufbauen, somit stimmt diese aussage

Die gewählten Parameter, wie beispielsweise der Steigungsfaktor $k$ oder der Wendepunkt $n_0$, dienen in der aktuellen Implementierung als Referenzwerte. Eine finale Optimierung dieser Hyperparameter erfolgt in zukünftigen Arbeiten durch eine empirische Evaluierung der Aggregationsergebnisse auf Basis realer Datensätze (siehe Kapitel~\ref{chapter:diskussion_und_forschungsluecke} (Diskussion und Forschungslücken)).

\begin{itemize}
    \item \textbf{Erzeugung fester Merkmalsvektoren:}
    Die Anforderung der Invarianz gegenüber variablen Eingabekardinalitäten wird durch die Funktion \texttt{aggregate\_weighted\_reports} (Listing~\ref{lst:aggregation} (\texttt{aggregation.py})) realisiert. In Anlehnung an den Attention-Mechanismus ermöglicht dieser Ansatz die Verarbeitung einer beliebigen Anzahl von Eingangsgrößen, während gleichzeitig ein Merkmalsvektor mit invarianter Dimension (Fixed-Size Input) extrahiert wird \cite{velickovic2018gat}. Dies stellt die notwendige Kompatibilität für nachgelagerte ML-Modelle sicher.

    \item \textbf{Relevanzbasierte Gewichtung:}
    Die qualitative Differenzierung der CTI-Reports erfolgt in der Funktion\texttt{count\_attack\_ongoing\_per\_report} (Listing~\ref{lst:attackCount} (\texttt{attack\_count.py})). Diese extrahiert die Anzahl der Einzeldetektionen innerhalb eines CTI-Reports und überführt diese als quantitative Basis in den Gewichtungsprozess. Damit wird sichergestellt, dass die Aggregation die Informationsdichte der einzelnen Quellen berücksichtigt.

    \item \textbf{Asymptotisches Sättigungsverhalten:}
    Um eine Verzerrung der Ergebnisse durch statistische Ausreißer oder extrem hohe Detektionszahlen zu verhindern, wurde ein asymptotisches Sättigungsverhalten implementiert. Die Funktionen \texttt{add\_report\_weight} und \texttt{add\_report\_weight\_sigmoid} (Listing~\ref{lst:aggregation} (\texttt{aggregation.py})) sind so konzipiert, dass sie gegen einen definierten Grenzwert konvergieren (siehe Kapitel~\ref{sec:auswahl_der_gewichtungsfunktion} (Auswahl der Gewichtungsfunktion)). Während die Utility-Theory-Funktion einen streng konkaven Verlauf aufweist, nutzt die Sigmoid-Funktion eine nichtlineare S-Kurve zur Stabilisierung der Gewichte (siehe Abbildung~\ref{fig:vergleich_gewichtungsfunktionen}), ähnlich der in Graph Attention Networks (GATs) verwendeten Sättigungs- oder Aktivierungsmechanismen \cite{velickovic2018gat, perezmaldonado2025sigmoid}. Eine vergleichende Evaluation beider mathematischen Modelle hinsichtlich ihrer Leistungsfähigkeit im operativen Betrieb ist Bestandteil zukünftiger Forschung (siehe Kapitel~\ref{chapter:diskussion_und_forschungsluecke} (Diskussion und Forschungslücken)).
\end{itemize}